% ===================================================================
%  TAULA N° 7 — Lo Vilatge a l'ivèrn. — Lo Mas a la prima.
%  Traduction 2026 d'apres texto/io, controlee sur texto/fr et
%  texto/ca. Consigne : voir texto/oc/10-taula-01.tex.
% ===================================================================

%%K t07-noto-1 noto
\VUnotes{143.2mm}{%
(*) Pel detalh d'un repais, vejatz las taulas 6 e 13.%
}

%%K t07-apar-1 apar
\VUsaut{4.0mm}
\VUcentre{13.2pt}{90}{SEGONDA SÈRIA}
\VUsaut{3.0mm}
\VUfilet{16mm}
\VUsaut{5.6mm}
\VUtitre{15.0pt}{60}{TAULA N° 7}
\VUsaut{3.0mm}

%%K t07-tit-1 sub
\VUcentre{12.0pt}{0}{\textit{Primièra scèna.}}
\VUsaut{3.0mm}

%%K t07-tit-2 sub
\VUpk{10.2pt}{40}{Lo vilatge a l'ivèrn.}
\VUsaut{4.0mm}

%%K t07-c1-01-1 p
1. --- Pendent las vacanças del Cap d'An acompanhèri mon paire a Vilanòva, un
vilatget ont aviá qualques afars a tractar.

%%K t07-c1-02-1 p
2. --- Prenguèrem la \VUgras{diligéncia} \textsuperscript{(11)} que fa lo servici
entre la vila e aquel vilatge; mas coma fasiá fòrça freg, aquel viatge dins un
vehicul tan pauc comòde foguèt pas brica agradiu.

%%K t07-c1-03-1 p
3. --- Per aquò sentiguèrem un veritable plaser quand vegèrem que lo
\VUgras{cochièr} arrestava sa veitura sus la plaça, davant l'\VUgras{ostalariá}
\textsuperscript{(2)} de l'\textit{Ors Blanc}. L'\VUgras{ostalièr}
\textsuperscript{(4)} nos esperava sul lindal de son ostal, en fumar
tranquillament sa \VUgras{pipa}.

%%K t07-c1-03-2 p
Nos convidèt a dintrar e me faguèri pas pregar per m'installar davant lo fuòc de
sarments que beluguejava dins lo fogal. Alara me trufava de l'aspre \VUgras{vent
del nòrd} que fasiá grinçar la vièlha \VUgras{ensenha} \textsuperscript{(3)} e
emportava en negres revolums lo \VUgras{fum} \textsuperscript{(1)} que sortiá de
la chiminèa.

%%K t07-c1-04-1 p
4. --- Èra l'ora de dinnar. Sens esperar mai, faguèrem onor al dinnar simple mas
saborós que nos serviguèron (*).

%%K t07-tit-3 sub
\VUpk{10.2pt}{40}{Qualques visitas.}
\VUsaut{3.0mm}

%%K t07-c1-05-1 p
5. --- Aprèp lo dinnar acompanhèri mon paire, qu'es agent d'assegurança, en çò de
sos clients. D'en primièr visitèrem lo \VUgras{fabre} \textsuperscript{(5)}, que
demòra contra l'ostalariá. Èra a ferrar un caval. Admirèri la fòrça de son
ajudaire, que tenia fermament contra son davantal de cuèr l'\VUgras{ongla} del
caval mentre que lo fabre fixava la \VUgras{ferradura} \textsuperscript{(6)} a
fòrts còps de martèl.

%%K t07-c1-06-1 p
6. --- Puèi anèrem en çò del \VUgras{sabatièr} \textsuperscript{(7)} del vilatge,
l'oncle Jaume. Encara qu'aguèsse per ensenha una bòta \VUgras{magnifica}, se
contentava de pausar de solas a un parelh de sabatas vièlhas e traucadas.

%%K t07-c1-06-2 p
Lo brave òme nos diguèt en risent: «A la vila èri “sabatièr” e aviái pas de
clients. Aicí soi “petaçaire” e i a de trabalh de sobras».

%%K t07-c1-07-1 p
7. --- Nos parlèt de son vesin lo \VUgras{barbièr} \textsuperscript{(8)}, que sa
practica es pas contenta. Sos \VUgras{rasors} son totjorn embrecats e sas
\VUgras{cisèlas} talhan pas jamai. Mas coma es lo sol barbièr del vilatge, i a pas
autre remèdi que se daissar rasar e talhar los pels per el. En mai es un òme avar
e antipatic.

%%K t07-c1-08-1 p
8. --- Per bonur, los autres \VUgras{vesins} li se semblan pas. Lo
\VUgras{carretièr} \textsuperscript{(9)}, per exemple, es un òme excellent,
trabalhador e servicial. Fa gaug de li portar una veitura a adobar. Son trabalh es
totjorn plan acabat.

%%K t07-c1-09-1 p
9. --- L'oncle Jaume se planhiá pas tanpauc del \VUgras{selièr}
\textsuperscript{(10)}, qu'ajuda de còps a cordurar los \VUgras{arnescs} quand lo
trabalh pressa.

%%K t07-c1-09-2 p
Mon paire li demandèt de sa familha. «Totes van plan. Ma felena es a
l'\VUgras{escòla} \textsuperscript{(12)}. Sa \VUgras{mestressa}
\textsuperscript{(13)} n'es fòrça contenta; es bona \VUgras{escolana}
\textsuperscript{(14)}. Es pas coma aquel bandit de mon filh; pensa pas qu'a
jogar. Lo \VUgras{mèstre} \textsuperscript{(15)} arriba pas a li ensenhar res».

%%K t07-tit-4 sub
\VUsaut{3.0mm}
\VUpk{10.2pt}{40}{Charradissas del vilatge.}
\VUsaut{3.0mm}

%%K t07-c1-10-1 p
10. --- Lo brave òme nos contèt puèi las novèlas del vilatge. Anavan bastir una
escòla novèla, qu'aquela qu'èra installada dins la \VUgras{comuna} èra venguda
tròp pichona. Autrament èra aisit, que sufisiá de crompar una partida del jardin
del \VUgras{molinièr}. Aquò li anariá fòrça plan al paure òme. Dempuèi que los
freges son arribats, las \VUgras{mòlas} del \VUgras{molin} \textsuperscript{(16)}
podián pas mai mòlre lo blat, que la \VUgras{ròda} \textsuperscript{(17)} èra
presa dins lo \VUgras{glaç} \textsuperscript{(25)} del \VUgras{riu}
\textsuperscript{(23)}.

%%K t07-c1-11-1 p
11. --- «E en parlar de bastissas, avètz vista nòstra novèla \VUgras{glèisa}
\textsuperscript{(18)}? Es fòrça pintoresca jos son mantèl de nèu. Los
\VUgras{corbasses} \textsuperscript{(19)}, que reconeisson pas mai son vièlh
cloquièr, se pausan tristament sul grand arbre del \VUgras{cementèri}
\textsuperscript{(20)}».

%%K t07-c1-12-1 p
12. --- L'idèa del cementèri remembrèt al sabatièr las personas que mon paire
aviá conegudas e qu'èran mòrtas l'an passat. «Quantas \VUgras{tombas}
\textsuperscript{(21)}, mon brave sénher, se son apondudas a las autras jos lo
\VUgras{ciprès} \textsuperscript{(22)}! La mòrt a emportat nòstre vièlh notari.
Tot lo vilatge assistiguèt a son \VUgras{enterrament}».

%%K t07-c1-13-1 p
13. --- «Vaquí son successor, a patinar sul riu. Venguèt i a una estona per que
li adobèsse la correja de son \VUgras{patin} \textsuperscript{(26)}, que s'èra
copada».

%%K t07-c1-14-1 p
14. --- «Vaquí tanben, a la riba del riu, lo novèl \VUgras{garda-campèstre}
\textsuperscript{(27)}. Es un ancian gendarma e la terror dels piòts del vilatge.
Fugisson tanlèu que veson sas \VUgras{gèstras} \textsuperscript{(29)} e son
\VUgras{quepí} \textsuperscript{(28)}».

%%K t07-c1-15-1 p
15. --- «A! Vaquí l'oncle Josèp, que pòrta un \VUgras{carri de fais de lenha}
\textsuperscript{(30)} pel fornièr. --- Es verai, diguèt mon paire, reconeissi son
atalatge de \VUgras{muòlas} \textsuperscript{(31)}. Me cal parlar amb el e
profitarai de l'escasença. A reveire, oncle Jaume».

%%K t07-c1-16-1 p
16. --- Just quand sortissiam, los enfants del vilatge daissavan l'escòla e se
lançavan en corrent sul \VUgras{lisador} \textsuperscript{(33)}, amb lo risc de
tombar e de se copar un membre dins la \VUgras{tombada} \textsuperscript{(34)}. Un
d'eles prenguèt sa \VUgras{luja} \textsuperscript{(32)} en çò del carretièr, i
assetèt un de sos companhs e partiguèt en corrent.

%%K t07-c1-17-1 p
17. --- D'autres se precipitèron cap a un \VUgras{molon de nèu}
\textsuperscript{(37)} contra lo \VUgras{pont} \textsuperscript{(40)} e
comencèron de bombardar lo \VUgras{monard de nèu} \textsuperscript{(39)} qu'avián
fach lo jorn d'abans e al qual avián pausat un capèl, una pipa e una cartabèla
d'escòla en bandolièra.

%%K t07-c1-18-1 p
18. --- Se soscavan pauc del vent gelat e del freg. La màger part avián pas
quitament de \VUgras{escarpa} \textsuperscript{(35)} e, per s'escalfar aprèp la
batèsta de \VUgras{bòlas de nèu} \textsuperscript{(38)}, lor sufisiá un ponhat de
\VUgras{castanhas} \textsuperscript{(42)} fòrça caudas crompadas a la tanta
Jaumeta, la \VUgras{venderèla}, que tremòla de freg jos son \VUgras{chale}
\textsuperscript{(43)} tròp prim.

%%K t07-tit-5 sub
\VUsaut{3.0mm}
\VUcentre{12.0pt}{0}{\textit{Segonda scèna.}}
\VUsaut{3.0mm}

%%K t07-tit-6 sub
\VUpk{10.2pt}{40}{Lo mas a la prima.}
\VUsaut{4.0mm}

%%K t07-c2-01-1 p
1. --- Malgrat totes los plasers de l'ivèrn, amb una jòia prigonda saludam lo
retorn de la \VUgras{prima}.

%%K t07-c2-01-2 p
Quin tablèu tan alègre es la cort d'un mas, al calabrun, al mes de mai!

%%K t07-c2-02-1 p
2. --- A l'orizont, lo \VUgras{solelh} \textsuperscript{(1)} se colca a plan
plan, en negant la \VUgras{campanha} \textsuperscript{(3)} de sos \VUgras{rais}
\textsuperscript{(2)} de porpra. Lo diligent \VUgras{laurador}
\textsuperscript{(6)} se despacha de laurar son \VUgras{camp}
\textsuperscript{(4)}.

%%K t07-c2-02-2 p
Amb son \VUgras{araire} \textsuperscript{(5)}, tirat per un vigorós \VUgras{caval}
\textsuperscript{(7)}, traça lo \VUgras{rèc} \textsuperscript{(8)} ont lo
\VUgras{semenaire} \textsuperscript{(9)} getarà a ponhats la \VUgras{grana} que
donarà las espigas dauradas.

%%K t07-c2-03-1 p
3. --- Los \VUgras{segaires} \textsuperscript{(11)}, amb lors dalhas, an talhat
l'èrba del prat, los fenaires an secat lo \VUgras{fen} \textsuperscript{(10)} en o
virant amb las \VUgras{forcas} \textsuperscript{(12)} e los rastèls, e vaquí que
tòrnan.

%%K t07-c2-03-2 p
Los uns van davant lo carri pesuc tirat per de buòus; pòrtan lor esplech sus
l'espatla. Los autres, mai pigres, se son installats comodament sul fen odorós.

%%K t07-c2-04-1 p
4. --- En passar salúdan amistosament lo \VUgras{jardinièr}
\textsuperscript{(16)}, ocupat dins lo \VUgras{vergièr} \textsuperscript{(13)} a
poar los \VUgras{arbres fruchièrs} e a entar los \VUgras{perièrs}
\textsuperscript{(14)}. Los \VUgras{prunièrs} \textsuperscript{(15)} son en flor
e, dins l'\VUgras{òrt} \textsuperscript{(17)} plan fumat, las legumas, cauls e
ensaladas son ja bèlas.

%%K t07-c2-04-2 p
Sul muret, un bèl \VUgras{pavon} \textsuperscript{(18)} espandís la coa per que
s'admiren sas magnificas plumas.

%%K t07-tit-7 sub
\VUpk{10.2pt}{40}{La cort del mas.}
\VUsaut{3.0mm}

%%K t07-c2-05-1 p
5. --- Las pòrtas de l'\VUgras{estable} \textsuperscript{(19)} son dobèrtas e se
veson la manjadoira e lo \VUgras{jaç de palha} \textsuperscript{(23)} de l'\VUgras{èga}
\textsuperscript{(21)} que lo \VUgras{varlet d'estable} \textsuperscript{(20)} ven
de desatalar. Lo \VUgras{poutre} \textsuperscript{(22)}, tan graciós amb sas
cambas longas, seguís sa maire en sautant.

%%K t07-c2-06-1 p
6. --- Contra lo \VUgras{potz} \textsuperscript{(29)}, las \VUgras{vacas}
\textsuperscript{(25)} van beure a la \VUgras{pila} \textsuperscript{(26)} de fusta
que lor sèrv d'abeurador. Lo \VUgras{vedèl} \textsuperscript{(24)} ne profita per
tetar, e lo \VUgras{buòu} \textsuperscript{(27)}, en sentir la bona odor del
fenièr, bramís alègrament e auça amb orgulh lo cap amb sas poderosas
\VUgras{banas} \textsuperscript{(28)}.

%%K t07-c2-07-1 p
7. --- Totes trabalhan. La \VUgras{serviciala} \textsuperscript{(32)} ten a la man
la \VUgras{còrda} \textsuperscript{(31)} del potz. Tira d'aiga amb un
\VUgras{selh} \textsuperscript{(30)} per abeurar lo bestial. Contra lo
\VUgras{granièr} \textsuperscript{(33)}, un òme rastèla los brins de palha, e
davant la pòrta de l'\VUgras{ostal} \textsuperscript{(34)}, una vièlha
\VUgras{filaira} \textsuperscript{(35)}, amb sa colonha e son \VUgras{fus}, fila lo
\VUgras{lin} o lo \VUgras{cambe} coma dins los temps passats.

%%K t07-tit-8 sub
\VUsaut{3.0mm}
\VUpk{10.2pt}{40}{Lo mainatgièr.}
\VUsaut{3.0mm}

%%K t07-c2-08-1 p
8. --- Lo \VUgras{mainatgièr} \textsuperscript{(36)} dirigís tot aqueste trabalh e
dona d'òrdres a sos varlets. Manda recampar l'\VUgras{èrsa}
\textsuperscript{(40)} e lo \VUgras{pic} \textsuperscript{(39)} que son estats
oblidats contra la \VUgras{cabana} \textsuperscript{(38)} del \VUgras{gos}
\textsuperscript{(37)}.

%%K t07-c2-09-1 p
9. --- Sos crits espaurugan una \VUgras{colomba} \textsuperscript{(41)}, que fugís
a tot vòl cap al \VUgras{colombièr} \textsuperscript{(42)}, e las \VUgras{galinas}
\textsuperscript{(43)}, que se despachan de tornar al \VUgras{galinièr}
\textsuperscript{(44)}.

%%K t07-c2-10-1 p
10. --- Davant lo \VUgras{jaç de las fedas} \textsuperscript{(48)}, vaicí lo vièlh
\VUgras{pastre} \textsuperscript{(45)} que torna amb son \VUgras{tropèl}
\textsuperscript{(49)}. Amb son \VUgras{gaiòl} \textsuperscript{(46)}, que lo
daissa pas jamai, coma tanpauc sa \VUgras{blòda} \textsuperscript{(47)}
descolorida, mena al jaç \VUgras{marrans} \textsuperscript{(50)},
\VUgras{fedas} \textsuperscript{(53)}, \VUgras{anhèls}
\textsuperscript{(54)}, \VUgras{cabras} \textsuperscript{(51)} e
\VUgras{cabrits} \textsuperscript{(52)}. Son fidèl \VUgras{gos}
\textsuperscript{(55)} Argos tampa la marcha e anima amb sos japadisses los
tardièrs.

%%K t07-c2-11-1 p
11. --- Tot aqueste bruch e aqueste movement destorban pas la calma de la bona
\VUgras{vaca lachièira} \textsuperscript{(56)}, que fa sonar sa \VUgras{esquila}
\textsuperscript{(57)} mentre que la mòlzon davant la pòrta de l'\VUgras{estable}
\textsuperscript{(58)}.

%%K t07-c2-12-1 p
12. --- Sembla comprene que deu donar son lach per que la \VUgras{mainatgièira}
\textsuperscript{(60)} pòsca far de \VUgras{burre} dins la \VUgras{burrièira}
\textsuperscript{(61)} o los excellents \VUgras{fromatges}
\textsuperscript{(64)} qu'escorron sus la pòst pausada sul \VUgras{cussòl}
\textsuperscript{(63)}.

%%K t07-c2-13-1 p
13. --- Tota la \VUgras{lachariá} \textsuperscript{(59)} la ten fòrça neta lo
\VUgras{varlet del mas} \textsuperscript{(65)}, que ven de lavar los \VUgras{dorcs
del lach} \textsuperscript{(62)} dins lo grand selh que pòrta a la man.

%%K t07-tit-9 sub
\VUsaut{3.0mm}
\VUpk{10.2pt}{40}{La cort de las volalhas.}
\VUsaut{3.0mm}

%%K t07-c2-14-1 p
14. --- Amb sos gròsses esclòps camina un pauc pesucament, e atal fa fugir lo
\VUgras{piòt} \textsuperscript{(68)}, las \VUgras{anedas} \textsuperscript{(66)},
los \VUgras{anets} \textsuperscript{(67)} e las \VUgras{aucas}
\textsuperscript{(69)}, que dobrisson bèl lo \VUgras{bèc} \textsuperscript{(70)} e
getan de crits inquiets.

%%K t07-c2-14-2 p
Lo \VUgras{gal} \textsuperscript{(71)}, mai bellicós, auça sa \VUgras{cresta}
\textsuperscript{(72)} roja, brandís las plumas de sa \VUgras{coa}
\textsuperscript{(73)} e geta un sonòr «quiquiriquí».

%%K t07-c2-15-1 p
15. --- Una \VUgras{cloca} \textsuperscript{(74)} pigòta sul \VUgras{femorièr}
\textsuperscript{(81)} amb sos \VUgras{poletons} \textsuperscript{(75)}. Lo
\VUgras{porcelet} \textsuperscript{(78)} cor cap a sa maire, l'enòrma
\VUgras{truèja} \textsuperscript{(77)} que vesèm contra la sot.

%%K t07-c2-16-1 p
16. --- Dins sas gàbias, los \VUgras{conilhs} \textsuperscript{(79)} manjan amb
avidesa l'\VUgras{èrba} \textsuperscript{(80)} tendra que lor pòrta la filha del
mainatgièr. Joaneta aima fòrça sos conilhs e, cada jorn, aprèp anar quèrre los
\VUgras{uòus} \textsuperscript{(76)} fresques al galinièr, ven visitar sos
amiguets.
\VUsaut{6.0mm}
\VUfilet{22mm}
